% !TeX program = xelatex
% Compile twice with XeLaTeX. All illustrations are native vector TikZ drawings.
\documentclass[12pt,a4paper,oneside]{report}
\XeTeXgenerateactualtext=1
\usepackage[margin=23mm,top=24mm,bottom=24mm,headheight=27pt]{geometry}
\usepackage{amsmath,amssymb,mathtools,bm}
\usepackage{booktabs,array,longtable,tabularx}
\usepackage[table]{xcolor}
\usepackage{tikz}
\usetikzlibrary{arrows.meta,calc,patterns}
\usepackage[most]{tcolorbox}
\usepackage{titlesec,needspace}
\usepackage{hyperref}
\usepackage{bookmark}
\usepackage{xepersian}
\settextfont[Scale=1.06,BoldFont=Amiri-Bold.ttf]{Amiri-Regular.ttf}
\setlatintextfont[BoldFont=texgyretermes-bold.otf,ItalicFont=texgyretermes-italic.otf,BoldItalicFont=texgyretermes-bolditalic.otf]{texgyretermes-regular.otf}
\definecolor{navy}{HTML}{163747}
\definecolor{teal}{HTML}{127C80}
\definecolor{gold}{HTML}{B18542}
\definecolor{paper}{HTML}{F3F7F7}
\definecolor{muted}{HTML}{5E7178}
\definecolor{alert}{HTML}{A45538}
\hypersetup{unicode=true,colorlinks=true,linkcolor=navy,urlcolor=teal,
 pdftitle={روش‌های آماری ـ جزوهٔ بازنویسی‌شده},pdfauthor={علیرضا جهانبخش و احسان شهبازی},pdfsubject={احتمال، برآورد نقطه‌ای، فاصلهٔ اطمینان و آزمون فرض}}
\setlength{\parindent}{0pt}
\setlength{\parskip}{5pt plus 1pt}
\linespread{1.12}
\raggedbottom
\setcounter{tocdepth}{1}
\setcounter{secnumdepth}{2}
\allowdisplaybreaks[2]
\emergencystretch=1em
% Page styles use explicit boxes to keep bidi direction groups balanced.
\makeatletter
\def\ps@lecture{%
  \def\@oddhead{\vbox{\hbox to\textwidth{\small
    \rl{روش‌های آماری}\hfil\rl{\leftmark}}%
    \kern4pt\hrule height .35pt}}%
  \let\@evenhead\@oddhead
  \def\@oddfoot{\hfil\thepage\hfil}%
  \let\@evenfoot\@oddfoot}
\makeatother
\pagestyle{lecture}
\renewcommand{\chaptermark}[1]{\markboth{#1}{}}
\titleformat{\chapter}[display]{\color{navy}\bfseries\Huge}{\large فصل \thechapter}{8pt}{}[\vspace{4pt}\color{gold}\titlerule]
\titlespacing*{\chapter}{0pt}{0pt}{18pt}
\titleformat{\section}{\Large\bfseries\color{navy}}{\thesection}{.6em}{}
\titleformat{\subsection}{\large\bfseries\color{teal}}{\thesubsection}{.6em}{}
\tcbset{enhanced,breakable,boxrule=.45pt,arc=1.5mm,
 left=3mm,right=3mm,top=2mm,bottom=2mm,before skip=9pt,after skip=9pt,
 fonttitle=\bfseries,coltitle=white,colback=paper,colframe=teal}
\newtcolorbox{definition}[1]{breakable=false,title={تعریف | #1},colframe=navy,colback=navy!3!white}
\newtcolorbox{result}[1]{breakable=false,title={نتیجهٔ کلیدی | #1},colframe=teal,colback=teal!4!white}
\newtcolorbox{important}[1]{breakable=false,title={دقت کنید | #1},colframe=gold,colback=gold!6!white}
\newcounter{ex}[chapter]
\newenvironment{example}[1]{\refstepcounter{ex}\begin{tcolorbox}[breakable=false,title={مثال \thechapter.\arabic{ex} | #1},colframe=muted,colback=white]}{\end{tcolorbox}}
\newcommand{\edit}[1]{\par\begingroup\small\color{alert}\textbf{تکمیل یا اصلاح علمی.} #1\par\endgroup}
\newcommand{\source}[1]{\par{\small\color{muted} مبنا: صفحه‌های #1 فایل دست‌نویس.}\par}
\newcommand{\E}{\mathbb E}
\newcommand{\Var}{\operatorname{Var}}
\newcommand{\Cov}{\operatorname{Cov}}
\newcommand{\MSE}{\operatorname{MSE}}
\newcommand{\Bias}{\operatorname{Bias}}
\newcommand{\SE}{\operatorname{SE}}
\newcommand{\iid}{\overset{\mathrm{iid}}{\sim}}
\newcommand{\N}{\mathcal N}
\newcommand{\R}{\mathbb R}
\newcommand{\Prob}{\mathbb P}
\newcommand{\sol}{\par\textbf{حل.} }
\newcommand{\figcap}[1]{\par\smallskip{\small\color{muted}#1}\par}

\hypersetup{pdftitle={برآورد نقطه‌ای و معیارهای کیفیت},pdfauthor={دکتر سکینه دهقان؛ علیرضا جهانبخش و احسان شهبازی}}
\begin{document}
\setcounter{chapter}{1}
\chapter{برآورد نقطه‌ای و معیارهای کیفیت}
\begin{center}{\small\color{muted} استاد درس: دکتر سکینه دهقان\quad | \quad گردآورندگان جزوه: علیرضا جهانبخش و احسان شهبازی}\end{center}

\source{۵ تا ۸}
\section{جامعه، نمونه، پارامتر و آماره}
\begin{definition}{مفاهیم پایهٔ استنباط}
\textbf{جامعهٔ آماری}: مجموعهٔ واحدهایی که می‌خواهیم دربارهٔ آن‌ها اطلاعات به دست آوریم.

\textbf{نمونه}: بخشی از واحدهای جامعه. نمونه‌گیری تصادفی به هر واحد، طبق طرح مشخص، امکان انتخاب می‌دهد؛ برابر بودن احتمال انتخاب، ویژگی نمونه‌گیری تصادفی ساده است.

\textbf{نمونهٔ تصادفی از یک توزیع}: متغیرهای $X_1,\ldots,X_n$ مستقل و هم‌توزیع‌اند؛ این فرض را با $\iid$ نشان می‌دهیم.

\textbf{پارامتر}: ویژگی عددی ثابت اما معمولاً نامعلومِ توزیع جامعه؛ مانند $\mu$ یا $\sigma^2$. مجموعهٔ مقادیر مجاز آن، فضای پارامتر $\Theta$ است.

\textbf{آماره}: تابعی از مشاهدات نمونه که به پارامتر نامعلوم وابسته نیست. آماره، پیش از مشاهدهٔ داده‌ها، متغیری تصادفی و دارای توزیع نمونه‌ای است.
\end{definition}
\textbf{برآوردگر} آماره‌ای مانند $T=T(X_1,\ldots,X_n)$ است که برای تخمین یک پارامتر به کار می‌رود. \textbf{برآورد} مقدار عددی $T(x_1,\ldots,x_n)$ پس از مشاهدهٔ داده‌هاست.
\begin{example}{توزیع نمونه‌ای میانگین}
جامعه‌ای شامل چهار مقدار $1,2,3,4$ داریم و با جایگذاری، دو مشاهده انتخاب می‌کنیم. هر زوج مرتب احتمال $1/16$ دارد. توزیع میانگین نمونه عبارت است از
\begin{center}\small
\begin{tabular}{c|ccccccc}
\toprule
$\bar x$&$1$&$1.5$&$2$&$2.5$&$3$&$3.5$&$4$\\
\midrule
$\Prob(\bar X=\bar x)$&$1/16$&$2/16$&$3/16$&$4/16$&$3/16$&$2/16$&$1/16$\\
\bottomrule
\end{tabular}\end{center}
برای نمونه، میانگین $2$ از زوج‌های $(1,3),(2,2),(3,1)$ حاصل می‌شود. میانگین جامعه $2.5$ و واریانس آن $1.25$ است؛ بنابراین $\E(\bar X)=2.5$ و $\Var(\bar X)=1.25/2=0.625$.
\end{example}
\section{نااریبی و اریبی}
\begin{definition}{برآوردگر نااریب}
برآوردگر $T$ برای $\theta$ نااریب است اگر برای تمام $\theta\in\Theta$،
\[\E_\theta(T)=\theta.\]
اریبی آن برابر است با
\[\Bias_\theta(T)=\E_\theta(T)-\theta.\]
\end{definition}
نااریبی به معنی برابری هر برآورد با پارامتر نیست؛ میانگین برآوردها در تکرار نمونه‌گیری برابر پارامتر است.
\begin{example}{نااریبی میانگین نمونه}
اگر $X_1,\ldots,X_n\iid\N(\theta,1)$، آنگاه
\[\E_\theta(\bar X)=\frac1n\sum_i\E_\theta(X_i)=\theta.\]
در واقع، این نتیجه برای هر نمونهٔ مستقل و هم‌توزیع با میانگین متناهی برقرار است و به نرمال بودن نیاز ندارد.
\end{example}
\begin{example}{چرا در واریانس نمونه بر $n-1$ تقسیم می‌کنیم؟}
جامعه میانگین $\mu$ و واریانس متناهی $\sigma^2$ دارد. تعریف کنید
\[\widehat\sigma_n^2=\frac1n\sum_{i=1}^n(X_i-\bar X)^2.\]
از همانی
\[\sum_i(X_i-\bar X)^2=\sum_i(X_i-\mu)^2-n(\bar X-\mu)^2\]
و $\E[(\bar X-\mu)^2]=\sigma^2/n$ نتیجه می‌شود
\[
\E\left[\sum_i(X_i-\bar X)^2\right]=(n-1)\sigma^2.
\]
پس
\[\E(\widehat\sigma_n^2)=\frac{n-1}{n}\sigma^2,\qquad
\Bias(\widehat\sigma_n^2)=-\frac{\sigma^2}{n}.\]
در مقابل، برای $n\ge2$،
\[\boxed{S^2=\frac1{n-1}\sum_i(X_i-\bar X)^2,\qquad \E(S^2)=\sigma^2.}\]
\end{example}
\section{خطای میانگین مربعات و کارایی نسبی}
\begin{result}{تجزیهٔ خطای میانگین مربعات}
برای برآورد $\theta$،
\[
\MSE_\theta(T)=\E_\theta[(T-\theta)^2]
=\Var_\theta(T)+[\Bias_\theta(T)]^2.
\]
در نتیجه، برای برآوردگر نااریب، $\MSE_\theta(T)=\Var_\theta(T)$ است.
\end{result}
\textbf{اثبات.} بنویسید $T-\theta=[T-\E(T)]+[\E(T)-\theta]$. پس از مجذور کردن و امید گرفتن، جملهٔ متقاطع به‌علت $\E[T-\E(T)]=0$ حذف می‌شود.

در این جزوه، $T_1$ را در مقدار پارامتر مورد نظر کاراتر از $T_2$ می‌نامیم اگر $\MSE_\theta(T_1)<\MSE_\theta(T_2)$. کارایی نسبی به صورت
\[
e(T_1,T_2)=\frac{\MSE_\theta(T_2)}{\MSE_\theta(T_1)}
\]
تعریف می‌شود؛ مقدار بیشتر از یک به سود $T_1$ است. برای دو برآوردگر نااریب، می‌توان نسبت واریانس‌ها را جایگزین کرد.
\begin{important}{مقایسه باید برای یک هدف مشترک باشد}
برای مقایسهٔ دو برآوردگرِ اریب، نسبت واریانس‌ها کافی نیست. هر دو خطا باید نسبت به \textbf{همان پارامتر هدف} محاسبه شوند. اگر ترتیب برتری با $\theta$ تغییر کند، هیچ‌یک در کل فضای پارامتر برتری یکنواخت ندارد.
\end{important}
\subsection{دو مثال از برآوردگرهای اریب}
\begin{example}{نمونهٔ دوجمله‌ای؛ خوانش حفظ‌شده از دست‌نویس}
فرض کنید $X_i\iid\operatorname{Bin}(k,\theta)$ و هدف برآورد $m=k\theta$ است. در متن مبنا دو آمارهٔ
\[A=\bar X,\qquad B=\frac{X_1-\bar X}{n}\]
برای مقایسه آمده‌اند. با $v=k\theta(1-\theta)$ داریم
\[
\E(A)=m,\quad\Var(A)=v/n,\qquad
\E(B)=0,\quad\Var(B)=\frac{v(n-1)}{n^3}.
\]
در محاسبهٔ واریانس $B$، کوواریانس $\Cov(X_1,\bar X)=v/n$ را لحاظ کرده‌ایم. بنابراین
\[
\MSE_m(A)=\frac vn,\qquad
\MSE_m(B)=\frac{v(n-1)}{n^3}+m^2.
\]
برای $0<\theta<1$،
\[
e(A,B)=\frac{n-1}{n^2}+\frac{nk\theta}{1-\theta}.
\]
\edit{صورت آمارهٔ دوم در دست‌نویس نامعمول است؛ خوانش «تفاضل تقسیم بر $n$» حفظ شده است. این آماره برای $k\theta$ نااریب نیست و مقایسهٔ صرفِ واریانس‌ها درست نخواهد بود.}
\end{example}
\begin{example}{نمونهٔ پواسون}
فرض کنید $X_i\iid\operatorname{Pois}(\theta)$ و $\theta>0$. دو برآوردگر $A=\bar X$ و $B=X_1+\bar X$ را برای $\theta$ مقایسه کنید.
\sol
\[
\E(A)=\theta,\quad\Var(A)=\theta/n,\qquad
\E(B)=2\theta,\quad\Bias_\theta(B)=\theta.
\]
به‌علت حضور $X_1$ در میانگین،
\[
\Var(B)=\theta+\frac\theta n+2\frac\theta n
=\theta\left(1+\frac3n\right).
\]
پس
\[
e(A,B)=\frac{\theta(1+3/n)+\theta^2}{\theta/n}
=n+3+n\theta>1.
\]
بنابراین $\bar X$ برای تمام $\theta>0$ کاراتر است.
\end{example}
\section{ترکیب خطی نااریب با کمترین واریانس}
\begin{example}{ضرایب نااریبی در یک نمونه}
اگر $X_i$ها میانگین $\mu$ داشته باشند، برای $T=\sum_i a_iX_i$ داریم
\[\E(T)=\mu\sum_i a_i.\]
پس برای نااریبی در تمام مقادیر $\mu$ لازم و کافی است $\sum_i a_i=1$. اگر علاوه بر آن، مشاهدات مستقل و واریانس‌های آن‌ها برابر $\sigma^2$ باشد،
\[\Var(T)=\sigma^2\sum_i a_i^2\ge\frac{\sigma^2}{n},\]
زیرا $\sum_i a_i^2\ge(\sum_i a_i)^2/n$. تساوی برای $a_i=1/n$ برقرار است؛ پس میانگین نمونه در این خانواده کمترین واریانس را دارد.
\end{example}
\begin{result}{وزن‌دهی معکوس واریانس}
اگر $T_1,T_2$ مستقل و نااریب برای $\theta$، با واریانس‌های مثبت $v_1,v_2$ باشند، ترکیب بهینهٔ نااریب آن‌ها عبارت است از
\[
T^*=\frac{v_2}{v_1+v_2}T_1+\frac{v_1}{v_1+v_2}T_2,
\qquad \Var(T^*)=\frac{v_1v_2}{v_1+v_2}.
\]
\end{result}
\textbf{اثبات.} شرط نااریبی $a_1+a_2=1$ است. تابع $a_1^2v_1+(1-a_1)^2v_2$ در $a_1=v_2/(v_1+v_2)$ کمینه می‌شود. اگر مانند یادداشت حاشیه‌ای متن مبنا $v_2=3v_1$ باشد، وزن‌ها به‌ترتیب $3/4$ و $1/4$ هستند.
\section{سازگاری}
\begin{definition}{نزدیک شدن احتمالی به پارامتر}
دنبالهٔ برآوردگرهای $T_n$ برای $\theta$ سازگار است اگر برای هر $\varepsilon>0$ و هر $\theta\in\Theta$،
\[
\Prob_\theta(|T_n-\theta|>\varepsilon)\longrightarrow0.
\]
این خاصیت را با $T_n\xrightarrow{P}\theta$ نشان می‌دهیم.
\end{definition}
با نامساوی مارکف روی $(T_n-\theta)^2$،
\[
\Prob_\theta(|T_n-\theta|>\varepsilon)
\le\frac{\MSE_\theta(T_n)}{\varepsilon^2}.
\]
پس $\MSE_\theta(T_n)\to0$ شرط کافی سازگاری است. به‌ویژه، اگر $T_n$ نااریب باشد و $\Var_\theta(T_n)\to0$، سازگار است. به‌طور کلی نیز کافی است اریبی و واریانس هر دو به صفر میل کنند.
\edit{میل کردن واریانس به صفر به‌تنهایی کافی نیست؛ مثلاً برآوردگر ثابت صفر برای پارامتر ناصفر، واریانس صفر دارد ولی سازگار نیست. شرط $\MSE\to0$ کافی است و ضرورت ندارد.}
\begin{example}{سازگاری میانگین}
اگر $X_i\iid\N(\theta,1)$، آنگاه $\MSE_\theta(\bar X_n)=1/n\to0$. بنابراین $\bar X_n\xrightarrow{P}\theta$.

همین استدلال برای $X_i\iid\operatorname{Pois}(\theta)$ با $\MSE=\theta/n$ و برای $X_i\iid\operatorname{Bin}(k,\theta)$، به‌منظور برآورد $\theta$ با $\bar X_n/k$، با $\MSE=\theta(1-\theta)/(nk)$ برقرار است.
\end{example}
\begin{important}{تمرین‌های ارجاع‌شده در متن مبنا}
در پایان این بخش، مثال‌های صفحه‌های ۱۰۵ تا ۱۱۲ و تمرین‌های صفحهٔ ۱۵۳، سؤال‌های ۱ تا ۵ کتاب معرفی شده‌اند. چون نام و متن کتاب در فایل موجود نیست، این ارجاع حفظ شده است و صورت سؤال‌های کتاب بازسازی نشده است.
\end{important}


\end{document}
